\section*{Self-Test Questions}

Test your knowledge with these representative problems. Answer all questions independently before reviewing the solutions in the answer key below.

\subsection*{Questions}
\begin{enumerate}[label=\textbf{Q\arabic*.}]
    \item A point charge $q_1 = +5.0\text{ nC}$ is located at the origin. A second point charge $q_2 = -3.0\text{ nC}$ is located on the y-axis at $y = 0.30\text{ m}$.
    \begin{enumerate}
        \item Compute the magnitude of the electrostatic force between the two charges.
        \item State whether this force is attractive or repulsive.
    \end{enumerate}
    
    \item A capacitor with capacitance $C_0$ is fully charged by a battery to a potential difference $V_0$. The capacitor is then disconnected from the battery. If the plate separation $d$ is subsequently doubled, determine what happens to:
    \begin{enumerate}
        \item The charge $Q$ on the plates.
        \item The capacitance $C$.
        \item The potential difference $V$ between the plates.
    \end{enumerate}
    
    \item A uniform electric field of magnitude $E = 3.0 \times 10^3\text{ V/m}$ points in the $+x$ direction. A test charge $q = +2.0\text{ }\mu\text{C}$ is moved from $x = 0.20\text{ m}$ to $x = 0.80\text{ m}$ along a straight path.
    \begin{enumerate}
        \item Find the potential difference $\Delta V = V_f - V_i$.
        \item Calculate the work done on the charge by the electric field.
    \end{enumerate}
\end{enumerate}

\vspace{15pt}

% ANSWER KEY BOX
\begin{tcolorbox}[
    enhanced,
    colback=lightbg,
    colframe=primary,
    arc=4pt,
    boxrule=1pt,
    left=15pt,
    right=15pt,
    top=12pt,
    bottom=12pt,
    title={\sffamily\bfseries\faLock\quad Answer Key \& Detailed Solutions},
    coltitle=white,
    colbacktitle=primary
]
\begin{itemize}[leftmargin=15pt]
    \item \textbf{A1:} 
    \begin{enumerate}
        \item Apply Coulomb's Law:
        \[ F_e = k_e \frac{|q_1 q_2|}{r^2} = \left(8.99 \times 10^9\right) \frac{(5.0 \times 10^{-9})(3.0 \times 10^{-9})}{(0.30)^2} \]
        \[ F_e = \left(8.99 \times 10^9\right) \frac{1.5 \times 10^{-17}}{0.09} = 1.50 \times 10^{-6}\text{ N} \]
        \item The force is \textbf{attractive} because the charges are of opposite signs.
    \end{enumerate}
    
    \item \textbf{A2:}
    \begin{enumerate}
        \item The charge remains constant ($Q' = Q_0$) because the capacitor is disconnected; there is no pathway for charge to leave.
        \item Capacitance is $C = \varepsilon_0 \frac{A}{d}$. Doubling $d$ halves the capacitance ($C' = \frac{1}{2} C_0$).
        \item Potential difference is $V = \frac{Q}{C}$. Since $Q$ is constant and $C$ is halved, the potential difference doubles ($V' = 2 V_0$).
    \end{enumerate}
    
    \item \textbf{A3:}
    \begin{enumerate}
        \item In a uniform field, $\Delta V = -E \Delta x = -(3.0 \times 10^3\text{ V/m})(0.80 - 0.20\text{ m}) = -1800\text{ V}$.
        \item The work done by the field is:
        \[ W_{field} = -q \Delta V = -(2.0 \times 10^{-6}\text{ C})(-1800\text{ V}) = +3.6 \times 10^{-3}\text{ J} = 3.6\text{ mJ} \]
    \end{enumerate}
\end{itemize}
\end{tcolorbox}
